\chapter{Matriz de Cabibbo--Kobayashi--Maskawa y violación de carga--paridad}
\label{chap:ckm}

La matriz de Cabibbo--Kobayashi--Maskawa (CKM) entrelaza la base de sabor de
los quarks con la base de autovectores de masa en tres generaciones. La
simetría de conjugación de carga y paridad se denota por CP\@. La
transformación racional que fija los tres ángulos no se introduce como una
tabla: se construye desde el sexteto de incidencia \(\pi\)--Witt y sus reglas
sectoriales. Después se evalúa sobre las coordenadas angulares
\[
 h_{\rm CKM}=\frac{\Cstar}{6},
 \qquad
 \Delta=\frac{180}{\pi}\Dfour.
\]
Las consecuencias de invertibilidad, unitariedad y no anulación del
invariante de Jarlskog se derivan entonces de esa construcción y de las
coordenadas internas.

\section{Genealogía interna de los coeficientes de cruce}
\label{sec:genealogia-n42-ckm-rev7}

El sexteto canónico de la incidencia de nivel \(N42\) es
\[
 (b,p,q,h_6,o_8,t)=(4,5,7,6,8,3).
\]
Sus tres reglas sectoriales son los vectores fila
\[
 R_{12}=(2,-p,bq),\qquad
 R_{23}=\left(0,q,-\frac{p^2}{2}\right),
\]
\[
 R_{13}=\left(1,-pb,-\frac{o_8-h_6}{t}\right).
\]

\begin{theorem}[Construcción sectorial de la transformación CKM]
\label{thm:construccion-n42-ckm-rev7}
Dadas las reglas anteriores y el sexteto \(N42\), la única transformación
que las tiene por filas es
\[
 \boxed{
 M_{\rm CKM}=
 \begin{pmatrix}
 2&-5&28\\
 0&7&-25/2\\
 1&-20&-2/3
 \end{pmatrix}.}
\]
Los coeficientes no dependen de un ajuste angular posterior.
\end{theorem}

\begin{proof}
La sustitución directa da
\[
 bq=4\cdot7=28,\qquad
 \frac{p^2}{2}=\frac{5^2}{2}=\frac{25}{2},
\]
\[
 pb=5\cdot4=20,qquad
 \frac{o_8-h_6}{t}=\frac{8-6}{3}=\frac23.
\]
Cada regla fija sus tres entradas; por tanto las tres filas y la matriz
completa quedan determinadas.
\end{proof}

La unicidad demostrada es relativa al sistema sectorial \(N42\) y a sus
seis invariantes. La realización física exige además entrelazar el espacio
de sabor quark con el espacio de masa; ese mapa no vuelve condicional la
construcción racional anterior ni se sustituye por el contraste PDG.

\section{Transformación racional de las coordenadas angulares}

Definimos
\[
 \begin{pmatrix}\theta_{12}\\\theta_{23}\\\theta_{13}\end{pmatrix}
 =M_{\rm CKM}
 \begin{pmatrix}A\\h_{\rm CKM}\\\Delta\end{pmatrix},
 \qquad
 M_{\rm CKM}=
 \begin{pmatrix}
 2&-5&28\\
 0&7&-25/2\\
 1&-20&-2/3
 \end{pmatrix},
\]
con
\[
 \delta_{\rm CP}=9A.
\]

\begin{theorem}[Invertibilidad del modelo angular CKM]
\[
 \det M_{\rm CKM}=-\frac{3857}{6}\ne0,
\]
y
\[
 M_{\rm CKM}^{-1}=
 \begin{pmatrix}
 1528/3857&3380/3857&801/3857\\
 75/3857&176/3857&-150/3857\\
 6/551&-30/551&-12/551
 \end{pmatrix}.
\]
Por tanto, los tres ángulos y las tres coordenadas generadoras contienen la
misma información lineal.
\end{theorem}
\begin{proof}
La expansión del determinante da \(-3857/6\). La multiplicación directa de
la matriz definida por la candidata a inversa produce la identidad.
\end{proof}

Se obtiene inmediatamente la relación lineal exacta
\[
\boxed{
 3857\,\delta_{\rm CP}
 =9(1528\theta_{12}+3380\theta_{23}+801\theta_{13}).}
\]
Esta ecuación pertenece a la parametrización estándar declarada en esta
sección. El ángulo \(\delta_{\rm CP}\) depende de esa parametrización; el
invariante de Jarlskog construido más adelante es la cantidad independiente
de las redefiniciones de fase.

\section{Matriz unitaria}

Los valores angulares se expresan en grados en las relaciones lineales y en
las tablas. En toda función trigonométrica o exponencial se emplea su
representante en radianes, \(\theta_{\rm rad}=\pi\theta/180\); para aligerar
la notación se conserva el mismo símbolo para dicho argumento.

Sea \(s_{ij}=\sin\theta_{ij}\), \(c_{ij}=\cos\theta_{ij}\). En la
parametrización estándar,
\[
V=
\begin{pmatrix}
c_{12}c_{13}&s_{12}c_{13}&s_{13}e^{-i\delta}\\
-s_{12}c_{23}-c_{12}s_{23}s_{13}e^{i\delta}&
c_{12}c_{23}-s_{12}s_{23}s_{13}e^{i\delta}&s_{23}c_{13}\\
s_{12}s_{23}-c_{12}c_{23}s_{13}e^{i\delta}&
-c_{12}s_{23}-s_{12}c_{23}s_{13}e^{i\delta}&c_{23}c_{13}
\end{pmatrix}.
\]

\begin{proposition}[Unitariedad]
Para cualesquiera ángulos reales y fase real, \(V^\dagger V=I_3\).
\end{proposition}
\begin{proof}
La matriz es el producto de tres rotaciones unitarias elementales, con la fase
insertada en la rotación \(13\). El producto de matrices unitarias es
unitario.
\end{proof}

La evaluación es
\[
 (\theta_{12},\theta_{23},\theta_{13},\delta_{\rm CP})
 =(13.003313589509,\;2.398056157332,
 0.213977382244,\;65.676173123554)^\circ.
\]
Los módulos son
\[
 |V|\simeq
 \begin{pmatrix}
 0.974350259&0.225005836&0.003734601\\
 0.224873110&0.973489279&0.041841465\\
 0.008582400&0.041121745&0.999117283
 \end{pmatrix}.
\]

\section{Invariante de Jarlskog}

El invariante independiente de base es
\[
 J=c_{12}c_{23}c_{13}^2s_{12}s_{23}s_{13}\sin\delta_{\rm CP}.
\]
En el punto anterior,
\[
 \boxed{J=3.1189723326632974\times10^{-5}\ne0.}
\]

\begin{theorem}[Violación CP en la matriz construida]
Si los cuatro parámetros del modelo se insertan en el sector de tres
generaciones, \(J\ne0\); por tanto la matriz no puede hacerse real mediante
redefiniciones de fase y la matriz contiene violación CP
\cite{Cabibbo1963,KobayashiMaskawa1973,Jarlskog1985}.
\end{theorem}
\begin{proof}
Los tres senos, los tres cosenos y \(\sin\delta_{\rm CP}\) son no nulos en la
evaluación. Su producto es el valor indicado, luego el invariante de Jarlskog
no desaparece.
\end{proof}

\section{Contraste con la referencia CKM}
\label{sec:comparacion-ckm-pdg-2026}

La \cref{tab:comparacion-angular-ckm-pdg-2026} confronta la evaluación
anterior, mantenida fija, con el ajuste global de tres generaciones publicado
en la revisión CKM de 2026 del Particle Data Group
\cite{PDGCKM2026}. El ajuste externo proporciona
\(\sin\theta_{12}\), \(\sin\theta_{13}\), \(\sin\theta_{23}\) y
\(\delta\) en radianes; los ángulos en grados se obtienen mediante la
conversión monótona de la parametrización estándar. Para una incertidumbre
asimétrica, el residuo normalizado emplea la rama correspondiente al signo de
la diferencia.

\begin{table}[htbp]
\centering
\caption{Parámetros angulares CKM e invariante de Jarlskog: evaluación
HMT--MD y ajuste global PDG 2026.}
\label{tab:comparacion-angular-ckm-pdg-2026}
\small
\begin{tabular}{@{}lrrr@{}}
\toprule
Parámetro & HMT--MD & PDG 2026 & Residuo normalizado\\
\midrule
\(\theta_{12}\) (grados) &
\(13.00331359\) & \(13.01287\pm0.03999\) & \(-0.24\)\\
\(\theta_{23}\) (grados) &
\(2.39805616\) & \(2.40082^{+0.04645}_{-0.03957}\) & \(-0.07\)\\
\(\theta_{13}\) (grados) &
\(0.21397738\) & \(0.215605^{+0.005042}_{-0.004756}\) & \(-0.34\)\\
\(\delta_{\rm CP}\) (grados) &
\(65.67617312\) & \(66.1193\pm1.4324\) & \(-0.31\)\\
\(10^5J\) &
\(3.11897233\) & \(3.16^{+0.13}_{-0.11}\) & \(-0.37\)\\
\bottomrule
\end{tabular}
\end{table}

Esta tabla es \textsc{reconocimiento externo}; no interviene en la
definición de \(M_{\rm CKM}\), en la evaluación de los cuatro parámetros ni
en la prueba de unitariedad. Los cinco residuos no constituyen cinco
contrastes independientes. El ajuste global es correlacionado y está
condicionado a la unitariedad de tres generaciones; sin su matriz de
covarianza, la tabla sólo ofrece diagnósticos marginales. El contraste
multivariado científicamente vinculante sigue siendo el
\(\chi^2\) definido más adelante.

La conclusión se refiere a CP\@. La simetría discreta CPT de rutas ya
construida actúa sobre otro objeto. Su realización como teorema cuántico de
campos se compara después con las hipótesis de localidad, covariancia de
Lorentz y positividad del espectro \cite{StreaterWightman2000}.

\section{Invertibilidad y comparación estadística}

La invertibilidad de \(M_{\rm CKM}\) demuestra que alterar uno de los tres
ángulos altera de manera determinada al menos una coordenada generadora. Para
comparar el modelo con un ajuste experimental debe usarse el vector de cuatro
parámetros independientes
\[
 y=(\sin\theta_{12},\sin\theta_{23},\sin\theta_{13},\delta_{\rm CP})
\]
y su covarianza completa:
\[
 \chi^2=(y-y_0)^{\!T}\Sigma^{-1}(y-y_0).
\]
Los nueve módulos de \(V\) no son nueve observables independientes: todos son
funciones de los cuatro parámetros. Esta reducción evita contar varias veces
la misma información.

La construcción \(N42\), la inversa racional, la unitariedad y el valor no
nulo de \(J\) forman una sola cadena: primero se determinan los coeficientes,
después se evalúan los ángulos y sólo al final se abre la referencia. La
identificación del entrelazador físico sabor--masa conserva su tipo propio;
no autoriza a rebajar \(M_{\rm CKM}\) a un dato ni a obtenerla desde los
valores observados.
